
\documentclass{article}
\usepackage{colortbl}
\usepackage{makecell}
\usepackage{multirow}
\usepackage{supertabular}

\begin{document}

\newcounter{utterance}

\centering \large Interaction Transcript for game `dond', experiment `semi\_en', episode 8 with Qwen3{-}32B{-}t0.0.
\vspace{24pt}

{ \footnotesize  \setcounter{utterance}{1}
\setlength{\tabcolsep}{0pt}
\begin{supertabular}{c@{$\;$}|p{.15\linewidth}@{}p{.15\linewidth}p{.15\linewidth}p{.15\linewidth}p{.15\linewidth}p{.15\linewidth}}
    \# & $\;$A & \multicolumn{4}{c}{Game Master} & $\;\:$B\\
    \hline

    \theutterance \stepcounter{utterance}  
    & & \multicolumn{4}{p{0.6\linewidth}}{
        \cellcolor[rgb]{0.9,0.9,0.9}{
            \makecell[{{p{\linewidth}}}]{
                \texttt{\tiny{[P1$\langle$GM]}}
                \texttt{You are playing a negotiation game in which you have to agree on how to divide a set of items among you and another player.} \\
\\ 
\texttt{Rules:} \\
\texttt{(a) You and the other player are given a set of items. Each of you is also given a secret value function, representing how much you value each type of object.} \\
\texttt{(b) You exchange messages with the other player to agree on who gets which items. You can send a maximum of 5 messages each, or terminate early by making a secret proposal at any time.} \\
\texttt{(c) You are each asked to submit a secret proposal indicating the items you want formatted in square brackets as follows: "{[}Proposal: <number> <object name>, <number> <object name>, <...>{]}"} \\
\texttt{(d) If your proposals are complementary, i.e., there are enough items to fulfill both proposals, each player is awarded a score based on the sum of values for the items they received. Otherwise, both of you get zero points.} \\
\texttt{(e) Your goal is to maximize the score you receive.} \\
\\ 
\texttt{Let us start.} \\
\\ 
\texttt{The set of available items is:} \\
\\ 
\texttt{3 shovels, 2 pictures, 1 shell} \\
\\ 
\texttt{Your secret value function is:} \\
\\ 
\texttt{shovel: 0, picture: 0, shell: 10} \\
\\ 
\texttt{IMPORTANT: Your messages, unless it is the secret proposal, are directly transmitted to the other player, so do not include any response to the rules or text announcing your message. To make a secret proposal, use the indicated format. Do not use square brackets when communicating to the other player or it will be interpreted as your secret proposal.} \\
            }
        }
    }
    & \\ \\

    \theutterance \stepcounter{utterance}  
    & \multicolumn{4}{p{0.6\linewidth}}{
        \cellcolor[rgb]{0.9,0.9,0.9}{
            \makecell[{{p{\linewidth}}}]{
                \texttt{\tiny{[P1$\rangle$GM]}}
                \texttt{<think>} \\
\texttt{Okay, let's see. I need to negotiate the division of items with another player. The available items are 3 shovels, 2 pictures, and 1 shell. My value function is shovel: 0, picture: 0, shell: 10. So, I only care about the shell. The other player's value function is secret, but I need to figure out how to get the shell for myself while ensuring their proposal doesn't conflict.} \\
\\ 
\texttt{First, since I value the shell so highly, I should aim to get that. But the other player might also want the shell. However, since there's only one shell, if both of us want it, our proposals would conflict, leading to zero points. So I need to find a way to get the shell without the other player needing it.} \\
\\ 
\texttt{Maybe the other player doesn't value the shell as much. Let's think: if I propose taking the shell, I need to make sure the other player's proposal doesn't require the shell. If the other player's proposal is for shovels or pictures, then we can both get what we want. Since I don't value shovels or pictures, maybe I can offer them to the other player in exchange for the shell.} \\
\\ 
\texttt{Wait, but how do I communicate that? I can send messages. Let me start by proposing that I take the shell and they take the rest. But I need to check if that's possible. If I take the shell, then the other player can take all the shovels and pictures. Since there are 3 shovels and 2 pictures, maybe they value those. If their value function is such that they want those, then our proposals would be complementary.} \\
\\ 
\texttt{So, my secret proposal should be {[}Proposal: 1 shell{]}. Then, I need to make sure the other player's proposal doesn't include the shell. If they propose taking some shovels and/or pictures, then we can both get our items. But if they also want the shell, then we both get zero. Therefore, I need to negotiate in a way that the other player gives up the shell.} \\
\\ 
\texttt{But how? Since I can't know their value function, maybe I can offer them the other items. For example, I can message them saying I want the shell and offer them all the shovels and pictures. If they accept, then their proposal would be for the shovels and pictures, and mine} \\
            }
        }
    }
    & & \\ \\

    \theutterance \stepcounter{utterance}  
    & & & \multicolumn{2}{p{0.3\linewidth}}{
        \cellcolor[rgb]{0.9,0.9,0.9}{
            \makecell[{{p{\linewidth}}}]{
                \texttt{\tiny{[GM$|$GM]}}
                \texttt{proposal} \\
            }
        }
    }
    & & \\ \\

    \theutterance \stepcounter{utterance}  
    & & & \multicolumn{2}{p{0.3\linewidth}}{
        \cellcolor[rgb]{0.9,0.9,0.9}{
            \makecell[{{p{\linewidth}}}]{
                \texttt{\tiny{[GM$|$GM]}}
                [0, 0, 1]
            }
        }
    }
    & & \\ \\

    \theutterance \stepcounter{utterance}  
    & & \multicolumn{4}{p{0.6\linewidth}}{
        \cellcolor[rgb]{0.9,0.9,0.9}{
            \makecell[{{p{\linewidth}}}]{
                \texttt{\tiny{[GM$\rangle$P2]}}
                \texttt{You are playing a negotiation game in which you have to agree on how to divide a set of items among you and another player.} \\
\\ 
\texttt{Rules:} \\
\texttt{(a) You and the other player are given a set of items. Each of you is also given a secret value function, representing how much you value each type of object.} \\
\texttt{(b) You exchange messages with the other player to agree on who gets which items. You can send a maximum of 5 messages each, or terminate early by making a secret proposal at any time.} \\
\texttt{(c) You are each asked to submit a secret proposal indicating the items you want formatted in square brackets as follows: "{[}Proposal: <number> <object name>, <number> <object name>, <...>{]}"} \\
\texttt{(d) If your proposals are complementary, i.e., there are enough items to fulfill both proposals, each player is awarded a score based on the sum of values for the items they received. Otherwise, both of you get zero points.} \\
\texttt{(e) Your goal is to maximize the score you receive.} \\
\\ 
\texttt{Let us start.} \\
\\ 
\texttt{The set of available items is:} \\
\\ 
\texttt{3 shovels, 2 pictures, 1 shell} \\
\\ 
\texttt{Your secret value function is:} \\
\\ 
\texttt{shovel: 1, picture: 1, shell: 5} \\
\\ 
\texttt{IMPORTANT: Your messages, unless it is the secret proposal, are directly transmitted to the other player, so do not include any response to the rules or text announcing your message. To make a secret proposal, use the indicated format. Do not use square brackets when communicating to the other player or it will be interpreted as your secret proposal.} \\
\\ 
\\ 
\texttt{The other player has submitted a secret proposal. It is now your turn to submit a proposal of your own.} \\
            }
        }
    }
    & \\ \\

    \theutterance \stepcounter{utterance}  
    & & & \multicolumn{4}{p{0.6\linewidth}}{
        \cellcolor[rgb]{0.9,0.9,0.9}{
            \makecell[{{p{\linewidth}}}]{
                \texttt{\tiny{[GM$\langle$P2]}}
                \texttt{<think>} \\
\texttt{Okay, let's see. The other player has already submitted a secret proposal, and now I need to come up with my own. The goal is to maximize my score based on my value function. Let me start by recalling the items available: 3 shovels, 2 pictures, 1 shell. My values are shovel: 1, picture: 1, shell: 5. The other player's proposal isn't visible to me, but since they made a secret proposal, I need to make sure that my proposal is complementary to theirs. If our proposals are complementary, we each get the items we asked for and score points. Otherwise, both get zero.} \\
\\ 
\texttt{Since I don't know what the other player asked for, I need to make a proposal that is likely to be complementary. The key here is to request items in such a way that even if the other player took some, there's still enough left for me. Let me think about possible scenarios.} \\
\\ 
\texttt{First, the shell is the highest value item for me. I should prioritize getting the shell if possible. But there's only one shell. If the other player also wants the shell, then we can't both get it. But since the other player's value function is secret, I don't know if they value the shell highly. However, since they made a secret proposal, maybe they took the shell. But I can't be sure. If I request the shell, and they also did, then our proposals aren't complementary, and we both get zero. That's risky. Alternatively, maybe I should request the shell and some lower{-}value items to ensure that even if they took some shovels or pictures, there's still enough for me.} \\
\\ 
\texttt{Wait, but my value for the shell is 5, which is much higher than the others. So if I can get the shell, that's 5 points. If I can't, maybe I can get multiple shovels or pictures. Let's see. The maximum I can get from shovels and pictures is 3*1 + 2*1 = 5, which is the same as the shell. So if I can get the shell, that's 5 points. If not, maybe I can get 3 shovels and 2 pictures for 5 points. But if I take the shell, I can't take any other items. Wait, no. The proposal is about what I want. If} \\
            }
        }
    }
     \\ \\

    \theutterance \stepcounter{utterance}  
    & & & \multicolumn{2}{p{0.3\linewidth}}{
        \cellcolor[rgb]{0.9,0.9,0.9}{
            \makecell[{{p{\linewidth}}}]{
                \texttt{\tiny{[GM$|$GM]}}
                \texttt{continue} \\
            }
        }
    }
    & & \\ \\

    \theutterance \stepcounter{utterance}  
    & & & \multicolumn{2}{p{0.3\linewidth}}{
        \cellcolor[rgb]{0.9,0.9,0.9}{
            \makecell[{{p{\linewidth}}}]{
                \texttt{\tiny{[GM$|$GM]}}
                \texttt{player was instructed to make a proposal} \\
            }
        }
    }
    & & \\ \\

\end{supertabular}
}

\end{document}
